\subsection{Ejemplo ilustrativo de una aplicación paralela en C}

La asimilación práctica de la biblioteca MPI se facilita significativamente mediante el análisis de un programa clásico de transmisión de mensajes, comúnmente denominado programa de saludos o \textit{greetings}. Este diseño de software elemental no solo permite verificar que la infraestructura de comunicación del clúster esté operando de manera correcta, sino que también introduce los fundamentos de la filosofía de paso de mensajes en sistemas distribuidos \cite{pacheco_2011}. A través de este modelo, el programador puede observar cómo un único código fuente es capaz de bifurcar su comportamiento lógico en tiempo de ejecución en función del identificador asignado a cada proceso independiente \cite{gropp_using_mpi_1999}.

El núcleo de este ejemplo práctico radica en la implementación del paradigma SPMD (\textit{Single Program, Multiple Data}), donde múltiples instancias ejecutan las mismas instrucciones sobre conjuntos de datos locales. Aunque MPI cuenta con más de un centenar de funciones, este código utiliza únicamente el subconjunto mínimo de seis rutinas fundamentales (\texttt{MPI\_Init}, \texttt{MPI\_Finalize}, \texttt{MPI\_Comm\_size}, \texttt{MPI\_Comm\_rank}, \texttt{MPI\_Send} y \texttt{MPI\_Recv}), las cuales son suficientes para escribir programas paralelos completamente funcionales \cite{grama_intro_parallel_2003, gropp_using_mpi_1999}. Para comprender a fondo el rol de cada una de las subrutinas invocadas en el programa, en la Tabla~\ref{tab:funciones_ejemplo} se detalla la semántica de las funciones esenciales de MPI utilizadas en la construcción de este binario. Luego, en la Figura~\ref{fig:codigo_greetings} se detalla la sintaxis completa en lenguaje C para llevar a cabo esta demostración fundamental.

\begin{table}[h]
	\centering
	\caption{Funciones fundamentales de la API de MPI utilizadas en el ejemplo. Elaboración propia con base en \cite{grama_intro_parallel_2003, gropp_using_mpi_1999, pacheco_2011, mpi_forum_2025}.}
	\label{tab:funciones_ejemplo}
	\begin{tabular}{|l|p{6.0cm}|p{5.0cm}|}
		\hline
		\textbf{Función MPI}                      & \textbf{Parámetros que recibe}                                                  & \textbf{Propósito y descripción operativa}                                                                                     \\ \hline
		\multirow{2}{*}{\texttt{MPI\_Init}}       & \texttt{int *argc}: Puntero al número de argumentos.                            & \multirow{2}{5.0cm}{Inicializa el entorno de ejecución de MPI y establece los contextos de comunicación iniciales.}            \\ \cline{2-2}
		                                          & \texttt{char ***argv}: Puntero al vector de argumentos.                         &                                                                                                                                \\ \hline
		\multirow{2}{*}{\texttt{MPI\_Comm\_size}} & \texttt{MPI\_Comm comm}: Comunicador a evaluar (ej. \texttt{MPI\_COMM\_WORLD}). & \multirow{2}{5.0cm}{Extrae el número total de procesos activos asociados al comunicador global especificado.}                  \\ \cline{2-2}
		                                          & \texttt{int *size}: Puntero entero donde se almacena el total de procesos.      &                                                                                                                                \\ \hline
		\multirow{2}{*}{\texttt{MPI\_Comm\_rank}} & \texttt{MPI\_Comm comm}: Comunicador a evaluar.                                 & \multirow{2}{5.0cm}{Obtiene el identificador entero único (rango) asignado al proceso que invoca la rutina.}                   \\ \cline{2-2}
		                                          & \texttt{int *rank}: Puntero entero donde se almacena el identificador único.    &                                                                                                                                \\ \hline
		\multirow{6}{*}{\texttt{MPI\_Send}}       & \texttt{const void *buf}: Dirección de memoria de los datos.                    & \multirow{6}{5.0cm}{Envía un mensaje de forma síncrona bloqueante hacia un destino utilizando una etiqueta específica.}        \\ \cline{2-2}
		                                          & \texttt{int count}: Cantidad de elementos en el búfer.                          &                                                                                                                                \\ \cline{2-2}
		                                          & \texttt{MPI\_Datatype datatype}: Tipo de dato MPI.                              &                                                                                                                                \\ \cline{2-2}
		                                          & \texttt{int dest}: Rango del proceso destino.                                   &                                                                                                                                \\ \cline{2-2}
		                                          & \texttt{int tag}: Etiqueta entera de control.                                   &                                                                                                                                \\ \cline{2-2}
		                                          & \texttt{MPI\_Comm comm}: Comunicador utilizado.                                 &                                                                                                                                \\ \hline
		\multirow{7}{*}{\texttt{MPI\_Recv}}       & \texttt{void *buf}: Búfer de memoria para almacenar datos.                      & \multirow{7}{5.0cm}{Recibe un flujo de datos desde un origen determinado, bloqueando el proceso hasta completar la operación.} \\ \cline{2-2}
		                                          & \texttt{int count}: Capacidad máxima de elementos.                              &                                                                                                                                \\ \cline{2-2}
		                                          & \texttt{MPI\_Datatype datatype}: Tipo de dato esperado.                         &                                                                                                                                \\ \cline{2-2}
		                                          & \texttt{int source}: Rango del proceso origen.                                  &                                                                                                                                \\ \cline{2-2}
		                                          & \texttt{int tag}: Etiqueta esperada.                                            &                                                                                                                                \\ \cline{2-2}
		                                          & \texttt{MPI\_Comm comm}: Comunicador utilizado.                                 &                                                                                                                                \\ \cline{2-2}
		                                          & \texttt{MPI\_Status *status}: Estructura con metadatos.                         &                                                                                                                                \\ \hline
		\texttt{MPI\_Finalize}                    & Ninguno (\texttt{void}).                                                        & Libera los recursos asignados por la biblioteca MPI y concluye formalmente el entorno paralelo.                                \\ \hline
	\end{tabular}
\end{table}

\begin{figure}[ht]
	\centering
	\includegraphics[width=0.85\textwidth]{Joni/greetings.png}
	\caption{Programa de saludos distribuido en MPI C. Elaboración propia con base en \cite{pacheco_2011}.}
	\label{fig:codigo_greetings}
\end{figure}


La estructura expuesta en la Figura~\ref{fig:codigo_greetings} revela el flujo de control dual que caracteriza a las aplicaciones distribuidas controladas por rangos numéricos. Mientras que los procesos con un rango superior a cero empaquetan y transmiten una cadena de caracteres hacia el nodo central, el proceso con rango cero se configura en un estado de escucha activa para recopilar y desplegar secuencialmente la información en la consola de salida estándar \cite{pacheco_2011}. El bloque condicional (\texttt{if... else}) demuestra a la perfección el modelo de ejecución SPMD, donde un único código fuente genera comportamientos distintos basándose en el identificador o rango (\texttt{my\_rank}) de cada proceso \cite{pacheco_2011}. Además, expone de manera clara el concepto de memoria distribuida: la variable \texttt{greeting} existe de forma privada e independiente en la memoria de cada proceso, por lo que la única manera de compartir su contenido con el proceso maestro es a través de la red usando un envío explícito \cite{pacheco_2011}.


En la Figura~\ref{fig:salida_greetings} se presenta la captura la impresión de salida del programa, donde se detalla el comando de invocación y la impresión secuencial de los mensajes de saludo transmitidos por cada uno de los rangos operacionales. El desarrollo de la salida impresa evidencia la ejecución del comando de lanzamiento \texttt{mpirun} configurado con un entorno de cuatro procesos en paralelo a través del parámetro estructural \texttt{-np 4}. Como resultado de esta instrucción, el proceso identificado con el rango cero toma inmediatamente el control de la consola estándar imprimiendo su propio saludo de forma local. De manera subsiguiente, las líneas de texto correspondientes a los rangos uno, dos y tres aparecen de forma ordenada debido al bucle síncrono de recepción soplado en el código maestro, el cual restringe la lectura consecutiva de las memorias buffer mediante llamadas deterministas a la primitiva \texttt{MPI\_Recv} \cite{open_mpi_2026}.


\begin{figure}[H]
	\centering
	\includegraphics[width=0.8\textwidth]{Joni/salida_greetings.png}
	\caption{Captura de pantalla de la terminal con la salida del programa \textit{greetings} en MPI. Elaboración propia.}
	\label{fig:salida_greetings}
\end{figure}


La correcta ejecución de este ejemplo de saludos demuestra la capacidad de las primitivas de comunicación punto a punto para establecer orden y sincronía en sistemas de memoria distribuida. El acoplamiento exacto entre las llamadas a \texttt{MPI\_Send} ejecutadas por los trabajadores y el bucle iterativo de \texttt{MPI\_Recv} coordinado por el nodo maestro asegura que los mensajes transiten por el tejido de red sin pérdidas ni colisiones de datos \cite{grama_intro_parallel_2003}. Este patrón de comunicación asimétrica constituye la base conceptual sobre la cual se cimientan algoritmos de descomposición de dominios más complejos en la computación científica actual \cite{mpi_forum_2025}.

A pesar de la aparente simplicidad matemática de este programa, es imperativo tomar precauciones rigurosas para evitar condiciones de bloqueo mutuo o \textit{deadlocks} durante la fase de escalamiento en supercomputadoras. Si el desarrollador altera el orden de las llamadas lógicas o si el proceso receptor intenta leer datos desde un rango inexistente, la aplicación completa quedará congelada indefinidamente en la memoria del sistema, desperdiciando valiosos recursos de procesamiento \cite{gropp_using_mpi_1999}. Por lo tanto, el diseño meticuloso de la topología lógica de los datos y el uso correcto de las estructuras de control representan las mejores directrices para la ingeniería de software de alto rendimiento \cite{open_mpi_2026}. Este programa es, en definitiva, el ejemplo por excelencia porque ilustra el ciclo de vida completo de una aplicación de MPI, enmarcando las comunicaciones estrictamente entre los llamados de inicio y fin, y comprender este patrón básico de comunicación y ramificación es el paso indispensable antes de avanzar a operaciones colectivas o arquitecturas de topologías virtuales \cite{grama_intro_parallel_2003, gropp_using_mpi_1999}.
